\clearpage\corrigesChapitre{9}{Synthétiser et identifier}
\corrigesSouples
\begin{blocLivre}{cor-9-1}\begin{corrige}{1}{L’origine d’une molécule}
\begin{questions}
\item Oui, si les entités sont identiques, il s’agit de la même espèce.
\item Non : un extrait naturel peut contenir beaucoup d’autres espèces ; un produit synthétique peut aussi présenter des impuretés.
\end{questions}
\end{corrige}\end{blocLivre}
\begin{blocLivre}{cor-9-2}\begin{corrige}{2}{Reconnaître le reflux}
\schema[corrige]{reflux}
\begin{questions}
\item Le ballon contient le mélange ; le réfrigérant vertical condense les vapeurs ; le chauffe-ballon apporte de l’énergie.
\item L’eau entre en bas et sort en haut pour remplir efficacement la gaine.
\end{questions}
\end{corrige}\end{blocLivre}
\begin{blocLivre}{cor-9-3}\begin{corrige}{3}{Le rôle du réfrigérant}
\begin{questions}
\item Il condense les vapeurs qui retournent dans le ballon et limite les pertes de matière.
\item Non : la transformation se déroule dans le mélange chauffé ; le réfrigérant agit surtout sur les vapeurs.
\end{questions}
\end{corrige}\end{blocLivre}
\begin{blocLivre}{cor-9-4}\begin{corrige}{4}{Une cuve de CCM}
\begin{questions}
\item Oui, le dépôt est au-dessus du niveau initial de l’éluant.
\item Les dépôts pourraient se dissoudre directement dans le bain et quitter la plaque.
\end{questions}
\end{corrige}\end{blocLivre}
\begin{blocLivre}{cor-9-5}\begin{corrige}{5}{Comparer deux produits}
\schema[corrige]{ccm-synthese}
\begin{questions}
\item Une tache de S est à la même hauteur que la tache de N dans les mêmes conditions.
\item Une seconde tache de S correspond à R : le réactif de départ est encore détecté.
\end{questions}
\end{corrige}\end{blocLivre}
\begin{blocLivre}{cor-9-6}\begin{corrige}{6}{Un montage à corriger}
\begin{questions}
\item Laisser une communication avec l’atmosphère et faire entrer l’eau en bas, sortir en haut.
\item Chauffer un système fermé peut provoquer une surpression.
\end{questions}
\end{corrige}\end{blocLivre}
\begin{blocLivre}{cor-9-7}\begin{corrige}{7}{Étapes d’un protocole}
\begin{questions}
\item Reflux : transformation ; filtration : isolement d’un solide ; CCM : contrôle.
\item Il réduit le risque lié à la chaleur et peut permettre la formation du solide recherché, selon le protocole.
\end{questions}
\end{corrige}\end{blocLivre}
\begin{blocLivre}{cor-9-8}\begin{corrige}{8}{Une seule tache}
\begin{questions}
\item Une seule espèce est détectée avec cette méthode et ces conditions.
\item Non : deux espèces peuvent migrer ensemble, ou une impureté ne pas être révélée.
\end{questions}
\end{corrige}\end{blocLivre}
\begin{blocLivre}{cor-9-9}\begin{corrige}{9}{Un test d’identité incomplet}
\begin{questions}
\item Non : différentes compositions peuvent donner des masses volumiques proches.
\item Une CCM dans les mêmes conditions et une température de changement d’état, si ces méthodes conviennent aux produits.
\end{questions}
\end{corrige}\end{blocLivre}
\begin{blocLivre}{cor-9-10}\begin{corrige}{10}{Légender utilement}
\schema[corrige]{reflux}
\begin{questions}
\item Elles n’identifient ni le réfrigérant, ni le ballon, ni les fonctions des éléments.
\item Le chauffage, le mélange, l’entrée et la sortie d’eau, ainsi que le retour des condensats et l’ouverture vers l’atmosphère.
\end{questions}
\end{corrige}\end{blocLivre}
\begin{blocLivre}{cor-9-11}\begin{corrige}{11}{Une comparaison équitable}
\begin{questions}
\item Non : la migration dépend de l’éluant et du support, ainsi que de la hauteur du front.
\item Déposer échantillon et référence sur une même plaque et développer ensemble.
\end{questions}
\end{corrige}\end{blocLivre}
\begin{blocLivre}{cor-9-12}\begin{corrige}{12}{Limiter l’impact}
\demarche{Examiner les critères fournis et préciser ceux qui manquent.}
\begin{questions}
\item B réduit ici le volume et la dangerosité annoncée du solvant.
\item Il faut aussi examiner l’énergie, les déchets, la quantité réellement obtenue et les conditions de sécurité.
\end{questions}
\retiens{Un choix de protocole demande plusieurs critères : produits, déchets, énergie et sécurité.}
\end{corrige}\end{blocLivre}

